\section{Technology Stack}

Each component below is given with a definition, a concrete example of its role
in this system, and the reason it was selected over alternatives.

\subsection{Model serving and inference}

\subsubsection{Ollama}
\begin{description}[leftmargin=0pt,style=nextline]
\item[Definition.] A local inference server that runs open-weight language
models from quantised GGUF files and exposes an HTTP API. It loads models on
demand and unloads them after an idle interval.
\item[Example.] \code{POST /api/chat} with \code{\{"model": "qwen3:14b",
"format": <json-schema>\}} returns output guaranteed to satisfy the schema.
\item[Why chosen.] Automatic load and unload gives VRAM management essentially
for free, which is decisive on a 16\,GB card. It ships CUDA~12.8 builds that
work on Blackwell today, and air-gapping is simply copying \code{\textasciitilde/.ollama}.
vLLM offers higher throughput but no automatic model swapping, which matters
more here than tokens per second.
\end{description}

\subsubsection{Open-weight models}
\begin{center}
\renewcommand{\arraystretch}{1.2}
\begin{tabularx}{\linewidth}{@{}l l r X@{}}
\toprule
\textbf{Role} & \textbf{Model} & \textbf{VRAM} & \textbf{Why this one} \\
\midrule
Reasoning, planning & \code{qwen3:14b} & 11.0\,GB & strongest planner that fits with headroom \\
Routing, extraction & \code{qwen3:4b} & 4.0\,GB & thousands of calls; must be cheap \\
Coding & \code{qwen2.5-coder:7b} & 6.0\,GB & specialised for code generation \\
Vision & \code{qwen2.5vl:7b} & 8.0\,GB & reads scans, handwriting, drawings \\
Embeddings & \code{bge-m3} & 1.2\,GB & dense \emph{and} sparse in one pass \\
\bottomrule
\end{tabularx}
\end{center}

\noindent Quantisation is 4-bit throughout. A \code{datacenter} profile
substitutes \code{gpt-oss:20b} where more memory is available, addressing the
problem statement's reference to 120B-class hardware.

\subsection{Retrieval and storage}

\subsubsection{K\`uzu --- embedded property graph database}
\begin{description}[leftmargin=0pt,style=nextline]
\item[Definition.] An in-process graph database with a Cypher-like query
language, analogous to what SQLite is for relational data.
\item[Example.] \code{MATCH (a:Equipment \{key:\$k\})-[r*1..3]-(b) RETURN b}
returns everything within three hops of an asset.
\item[Why chosen.] No server process and no listening port. On an air-gapped
box this is a security property, not merely a convenience.
\end{description}

\subsubsection{Qdrant (local mode) --- vector store}
\begin{description}[leftmargin=0pt,style=nextline]
\item[Definition.] A similarity search engine over embedding vectors; local
mode runs as an on-disk library rather than a service.
\item[Example.] A 1024-dimensional \code{bge-m3} embedding is stored per chunk
with its \code{doc\_id} and page; cosine search returns nearest passages.
\item[Why chosen.] Native support for combined dense and sparse vectors, and
the ability to run embedded.
\end{description}

\subsubsection{BM25 --- lexical ranking}
\begin{description}[leftmargin=0pt,style=nextline]
\item[Definition.] A term-frequency ranking function that scores documents by
exact word overlap, weighted by rarity.
\item[Example.] The query \code{P-101A} matches the literal token; an embedding
model may instead return semantically similar but wrong assets such as P-101B.
\item[Why chosen.] Exact identifier matching is a hard requirement in
industrial corpora, and it is precisely where dense retrieval is weakest.
\end{description}

\subsubsection{Reciprocal Rank Fusion}
\begin{description}[leftmargin=0pt,style=nextline]
\item[Definition.] A method of combining ranked lists using
$\mathrm{score}(d) = \sum_{i} 1/(k + \mathrm{rank}_i(d))$, where $k$ is a
smoothing constant, conventionally 60.
\item[Example.] A passage ranked 1st by BM25 and 5th by vector search scores
$1/61 + 1/65 = 0.0318$.
\item[Why chosen.] It is scale-free. Cosine similarities and BM25 scores live
on incomparable scales, and fusing by rank avoids ever calibrating them.
\end{description}

\subsubsection{Leiden community detection}
\begin{description}[leftmargin=0pt,style=nextline]
\item[Definition.] A graph clustering algorithm that partitions a network into
densely connected communities, improving on Louvain by guaranteeing
well-connected clusters.
\item[Example.] Applied to the indexed corpus it produced 37 communities, each
summarised once so that ``what are the recurring themes'' can be answered.
\item[Why chosen.] Thematic questions have no single source passage. Only
pre-computed cluster summaries can answer them.
\end{description}

\subsection{Document understanding}

\subsubsection{RapidOCR (ONNX Runtime)}
\begin{description}[leftmargin=0pt,style=nextline]
\item[Definition.] An optical character recognition pipeline --- text
detection, angle classification, recognition --- executed through ONNX Runtime
on the CPU.
\item[Example.] A skewed, noisy scan of an inspection report was processed in
2.4 seconds, recovering every asset tag correctly.
\item[Why chosen.] It requires no GPU and no PyTorch, so it leaves the entire
card for the language model, and it installs cleanly offline.
\end{description}

\subsubsection{Vision-language model}
\begin{description}[leftmargin=0pt,style=nextline]
\item[Definition.] A model accepting interleaved image and text input, able to
reason about layout, symbols and handwriting rather than only transcribe.
\item[Example.] Asked to list visible tags and the vibration reading in a
scanned report, it returned P-101A, TI-2043, E-204 and
``7.2\,mm/s RMS (alarm 4.5, trip 9.0)''.
\item[Why chosen.] OCR transcribes; it does not comprehend. A P\&ID requires
understanding of symbols and connectivity, not character recognition.
\end{description}

\subsection{Agent and tooling}

\subsubsection{Constrained decoding}
\begin{description}[leftmargin=0pt,style=nextline]
\item[Definition.] Restricting token generation so output provably conforms to
a supplied JSON schema.
\item[Example.] Entity extraction cannot emit a type outside the ontology
enumeration, because such tokens are masked during sampling.
\item[Why chosen.] It removes an entire class of parsing failure. Prompting a
model to ``reply in JSON'' and then parsing prose is unreliable at scale.
\end{description}

\subsubsection{Linux namespaces}
\begin{description}[leftmargin=0pt,style=nextline]
\item[Definition.] A kernel feature partitioning global resources so processes
see isolated views. A network namespace provides a private stack with only a
down loopback interface.
\item[Example.] \code{unshare -rn python main.py} runs generated code where
\code{socket.create\_connection(("1.1.1.1", 53))} raises
\code{OSError: Network is unreachable}.
\item[Why chosen.] It requires no root, and the guarantee is enforced by the
kernel rather than by application policy.
\end{description}

\subsubsection{SymPy}
\begin{description}[leftmargin=0pt,style=nextline]
\item[Definition.] A symbolic mathematics library performing exact algebraic
manipulation rather than floating-point approximation.
\item[Example.] Net positive suction head, evaluated symbolically then
substituted, returns $11.810537$ with every intermediate step recorded.
\item[Why chosen.] The problem statement requires ``calculations with steps
shown''. Language models perform arithmetic unreliably and produce no auditable
derivation; symbolic evaluation produces both.
\end{description}

\subsubsection{python-docx, python-pptx, openpyxl}
\begin{description}[leftmargin=0pt,style=nextline]
\item[Definition.] Libraries constructing genuine Office Open XML documents.
\item[Example.] An approval note is emitted with a title, numbered findings, a
three-column signature table and a provenance footer.
\item[Why chosen.] Filling templates from structured fields is far more
reliable than asking a model to emit document XML, and the result resembles the
organisation's own paperwork.
\end{description}

\subsection{Application and platform}

\begin{center}
\renewcommand{\arraystretch}{1.2}
\begin{tabularx}{\linewidth}{@{}l X@{}}
\toprule
\textbf{Component} & \textbf{Role and rationale} \\
\midrule
FastAPI + Uvicorn & Async HTTP API bound to \code{127.0.0.1}; no external interface \\
Pydantic & Typed settings and schema validation at every boundary \\
httpx & Single auditable HTTP transport, pinned to loopback by assertion \\
NetworkX + python-igraph & Graph construction and Leiden clustering \\
PyTorch (cu128) & CUDA runtime for Blackwell \code{sm\_120}; used by ONNX paths \\
nftables & Kernel firewall implementing default-deny egress \\
SQLite & Extraction cache and community summaries; zero configuration \\
\bottomrule
\end{tabularx}
\end{center}

\begin{keybox}{Verified platform}
Python 3.11.15 in a project-local virtualenv; PyTorch 2.11.0+cu128 confirmed on
compute capability $(12,0)$; NVIDIA RTX~5080 16\,GB; 32 CPU cores; 59\,GB system
RAM. Ninety-nine packages, all installed inside the virtualenv; nothing was
installed system-wide.
\end{keybox}
